\documentclass[11pt]{article}

% Change "review" to "final" to generate the final (sometimes called camera-ready) version.
% Change to "preprint" to generate a non-anonymous version with page numbers.
\usepackage[review]{acl}

% Standard package includes
\usepackage{times}
\usepackage{latexsym}

% For proper rendering and hyphenation of words containing Latin characters (including in bib files)
\usepackage[T1]{fontenc}
\usepackage{amsmath}
\usepackage{booktabs}
\usepackage{array}
\usepackage{tabularx}
\usepackage{ragged2e}
% For Vietnamese characters
% \usepackage[T5]{fontenc}
% See https://www.latex-project.org/help/documentation/encguide.pdf for other character sets

% This assumes your files are encoded as UTF8
\usepackage[utf8]{inputenc}

% This is not strictly necessary, and may be commented out,
% but it will improve the layout of the manuscript,
% and will typically save some space.
\usepackage{microtype}

% This is also not strictly necessary, and may be commented out.
% However, it will improve the aesthetics of text in
% the typewriter font.
\usepackage{inconsolata}

%Including images in your LaTeX document requires adding
%additional package(s)
\usepackage{graphicx}

\newcolumntype{Y}{>{\RaggedRight\arraybackslash}X}

% If the title and author information does not fit in the area allocated, uncomment the following
%
%\setlength\titlebox{<dim>}
%
% and set <dim> to something 5cm or larger.

\title{Instructions for *ACL Proceedings}

% Author information can be set in various styles:
% For several authors from the same institution:
% \author{Author 1 \and ... \and Author n \\
%         Address line \\ ... \\ Address line}
% if the names do not fit well on one line use
%         Author 1 \\ {\bf Author 2} \\ ... \\ {\bf Author n} \\
% For authors from different institutions:
% \author{Author 1 \\ Address line \\  ... \\ Address line
%         \And  ... \And
%         Author n \\ Address line \\ ... \\ Address line}
% To start a separate ``row'' of authors use \AND, as in
% \author{Author 1 \\ Address line \\  ... \\ Address line
%         \AND
%         Author 2 \\ Address line \\ ... \\ Address line \And
%         Author 3 \\ Address line \\ ... \\ Address line}

\author{First Author \\
  Affiliation / Address line 1 \\
  Affiliation / Address line 2 \\
  Affiliation / Address line 3 \\
  \texttt{email@domain} \\\And
  Second Author \\
  Affiliation / Address line 1 \\
  Affiliation / Address line 2 \\
  Affiliation / Address line 3 \\
  \texttt{email@domain} \\}

%\author{
%  \textbf{First Author\textsuperscript{1}},
%  \textbf{Second Author\textsuperscript{1,2}},
%  \textbf{Third T. Author\textsuperscript{1}},
%  \textbf{Fourth Author\textsuperscript{1}},
%\\
%  \textbf{Fifth Author\textsuperscript{1,2}},
%  \textbf{Sixth Author\textsuperscript{1}},
%  \textbf{Seventh Author\textsuperscript{1}},
%  \textbf{Eighth Author \textsuperscript{1,2,3,4}},
%\\
%  \textbf{Ninth Author\textsuperscript{1}},
%  \textbf{Tenth Author\textsuperscript{1}},
%  \textbf{Eleventh E. Author\textsuperscript{1,2,3,4,5}},
%  \textbf{Twelfth Author\textsuperscript{1}},
%\\
%  \textbf{Thirteenth Author\textsuperscript{3}},
%  \textbf{Fourteenth F. Author\textsuperscript{2,4}},
%  \textbf{Fifteenth Author\textsuperscript{1}},
%  \textbf{Sixteenth Author\textsuperscript{1}},
%\\
%  \textbf{Seventeenth S. Author\textsuperscript{4,5}},
%  \textbf{Eighteenth Author\textsuperscript{3,4}},
%  \textbf{Nineteenth N. Author\textsuperscript{2,5}},
%  \textbf{Twentieth Author\textsuperscript{1}}
%\\
%\\
%  \textsuperscript{1}Affiliation 1,
%  \textsuperscript{2}Affiliation 2,
%  \textsuperscript{3}Affiliation 3,
%  \textsuperscript{4}Affiliation 4,
%  \textsuperscript{5}Affiliation 5
%\\
%  \small{
%    \textbf{Correspondence:} \href{mailto:email@domain}{email@domain}
%  }
%}

\begin{document}
\maketitle
\begin{abstract}
Human--AI creative co-writing is increasingly popular for open ended text generation, yet most evaluations focus on final outputs rather than the interactional dynamics that produce the stories. We investigate narrative agency in turn-based co-writing across three settings: Human--Human, Human--AI, and AI--AI collaboration. Building on a previously introduced Human--AI corpus, we add Human--Human and AI--AI conditions under the same storytelling paradigm, enabling controlled comparison across pairings. We use an analytical framework that combines measures of affective alignment, semantic novelty, transience, and resonance to model how writers adapt to one another and shape subsequent narrative development. Our analysis shows that Human--AI co-writing constitutes a distinct asymmetric regime: LLMs accommodate human emotional tone and elaborate prior context, while humans introduce more novel and more persistent narrative material. In contrast, Human--Human and AI--AI interactions exhibit more symmetric patterns of adaptation. These findings suggest that LLMs function as adaptive amplifiers, both stabilizing and extending human-led narrative trajectories.
\end{abstract}

\section{Introduction}

Human-AI creative co-writing is becoming a common mode of open-ended text generation \citep{noy2023experimental,clark-smith-2021-choose,cheng2025inspiration}, as human creative writers, professional or otherwise, integrate large language models (LLMs) into their workflows. Human writers tend to treat these systems, trained on vast textual corpora, as collaborators capable of contributing to narrative production \citep{lee2022coauthor, wan2024felt}. As such interactions become widespread, it becomes important to understand how these models participate in the unfolding, turn-by-turn dynamics of collaborative storytelling. 

Turn-based co-writing produces a controlled setting for studying human-LLM interaction as a dynamic process. In this setting, writing unfolds as a sequence of interdependent contributions: each turn conditions the next turn and may shift the story's semantic direction, emotional trajectory and subsequent context. Consequently, turn-based narrative generation makes it possible to track how emotional tone, semantic novelty and narrative influence move between writers over time. While work on human-LLM turn-based dynamics is a growing field, such studies often focus on final outputs, isolated agent contributions, or aggregated metrics, leaving less attention to the fine-grained interactional dynamics that construct a story.

In this work, we move from single-condition analysis to a comparative framework of collaborative dynamics. We test whether narrative agency and emotional alignment differ across three co-writing settings: Human-AI (HA), Human-Human (HH) and LLM-LLM (AA). Specifically, we examine how different pairings allocate three core functions of storytelling: (1) The emotional alignment between agents; (2) The generation of novel narrative material; (3) The adaptation to the evolving context of the story.

We use a turn-level analytical framework that models co-writing as a dynamic process of interactions, designed for both agents to have theoretical equal autonomy in the writing process. By modeling contributions at the level of individual turns, the framework measures affective alignment, semantic novelty, and directional influence, allowing us to trace how (human and artificial) writers respond to prior context and shape subsequent narrative development. 

This comparative perspective allows us to ask \textbf{whether the asymmetries observed in human-AI interaction reflect general features of collaborative writing}, or specific properties of human-LLM narrative coordination. 

We use this framework to investigate the following nuances of asymmetries:

\begin{enumerate}
    \item \textbf{RQ1: Are dynamics of immediate emotional alignment symmetric or asymmetric across conditions?}
    \item \textbf{RQ2: Are agent-level differences in novel contribution and narrative influence coordination a general property of dyadic co-writing or intrinsic for human-LLM interactions?}
    \item \textbf{RQ3: Do narrative coordination in human-AI stories resemble human-human collaboration, LLM-LLM collaboration, or a distinct regime?}
\end{enumerate}

Our contributions are threefold. First, building on an existing Human--AI co-writing corpus, we introduce two new datasets of Human--Human and AI--AI co-written narratives under the same turn-based storytelling paradigm as the existing corpus. Second, we use a turn-level framework to expand methodologies for measuring affective alignment, measured as immediate valence adaptation of each agent, and discourse influence, measured as the relationship between novelty-transience-resonance. Finally, we use our novel to show that Human--AI co-writing forms a distinct asymmetric regime, unlike Human--Human or AI--AI.

\section{Related Works}

Human--AI co-writing is often motivated by the promise of hybrid intelligence: the possibility that human and artificial writers together may produce forms of creativity that neither could achieve alone \cite{cheng2025inspiration}.

Prior work on human–AI creativity has focused on evaluating the quality and productivity of generated outputs, finding mixed effects on creativity and ideation. 

A second line of work studies alignment and adaptation in dialogue, showing that human and artificial agents exhibit forms of linguistic and emotional coordination. 

Finally, computational approaches to narrative analysis have modeled sentiment trajectories and information flow in static texts. 

However, we lack a framework for analyzing how creative roles and narrative influence emerge during dynamic collaborative writing.

%\section{Framework and Definitions}

\section{Datasets}

% Ida and Charlotte territory
% Yb: I suggest a diagram of hh, ha, aa showing how they follow the same 10 interactions setup.

% Also a Table: dataset statistics: number of stories, participants and models, turns, mean words per turn and models used. 
\begin{table}[t]
\centering
\small
\begin{tabular}{lrrr}
\toprule
Condition & Stories & Turns & Mean words/turn \\
\midrule
HH & 38 & 342 & 26.13 \\
HA & 87 & 783 & 26.96 \\
AA & 40 & 360 & 28.24 \\
\bottomrule
\end{tabular}
\caption{Dataset statistics across co-writing conditions.}
\label{tab:dataset}
\end{table}

We use the same turn-based storytelling paradigm across three dyadic conditions, to compare collaborative dynamics holding a constant task structure.\footnote{We will release our dataset and analysis framework upon acceptance} To the best of our knowledge, this is the first dataset that enables direct comparison of Human-Human, Human-AI and LLM-LLM turn-based co-writing under a shared experiment paradigm. 

The following instruction was provided to all agents across conditions:
\begin{quote}
\textit{“You are an author taking part in a collaborative storytelling activity with another author. Together, you will create a story by taking turns adding to it. Your goal is to continue from where your partner has left off. If there is no story, please begin the story. You have 10 interactions to write the story. Your input may be slightly truncated by a random number of characters.”}
\end{quote}

The resulting datasets consists of the following four levels of analysis:

\begin{enumerate}
    \item \textbf{Token level (model level):} individual tokens $w$.
    \item \textbf{Turn level:} a single agent contribution at turn $t$, denoted $I_t$, which can be further specified as either agent 1's turn ($A^1_t$) or agent 2's turn ($A^2_t$).
    \item \textbf{Interaction level:} paired contributions at the same turn $t$, represented as either $(A^1_t, A^2_t)$ or the opposite direction $(A^2_{t-1}, A^1_t)$ .
    \item \textbf{Story level:} a sequence of 10 interactions for session $i$, denoted $S_i$.
\end{enumerate}


\subsection{Human-AI}
% cb: use either a citation from Arxiv or from ACL, depending on if F&B get's approved or not. For now I have used the Arxiv version.

The Human-AI data was collected by Fundal and Bizzoni \citeyearpar{fundal2026directionalalignmentnarrativeagency}. The condition was designed to isolate the turn-taking between the human and the AI in a collaborative storytelling exercise. There were X participants (male = X, female = X, other = X) with an age distribution of \textit{M} = X years (\textit{SD} = X). A link for participation was distributed through several online channels. The LLMs used were gpt-4.1-2025-04-14, claude-sonnet-4-5-20250929, Llama-3.3-70B-Instruct, and Qwen-2.5-72B-Instruct.


\subsection{Human-Human} 
The condition consisted of 76 participants (male = X, female = X, other = X) with an age distribution of \textit{M} = X years (\textit{SD} = X). All participants were university or graduate students at Aarhus University in Denmark, and participation was voluntary and uncompensated. All students were enrolled in either Linguistics, English, English International Business Communication, Cultural Data Science, or Cognitive Science. All degree programs are taught in English. Participants who had participated in the human-AI condition were excluded from the human-human condition. 
%ij: the number of participants is just the stories*2, but if you have a different number somehow, please change it! this might also be a great place to put more detailed demographics in (gender & age) but i don't have access to those.

Data collection was facilitated by student assistants and took place during a class or immediately after. The Human-Human condition used the same platform as the Human-AI condition. By accessing the platform, the participants were paired with another participant who was also queued for the writing session. They were not informed who their partner was, only that it was another human writer. The interface would allow one participant to write by showing a text-box, while the other person had no text-box to write in. Once the first participant had submitted their writing, the participant that they were paired up with would see the writing on their interface and have a text-box to continue the writing. Each participant was tasked with writing 10 sections for a total of 20 turns.

This condition resulted in 38 stories, with a total of 342 interactions after preprocessing and filtering. Each turn had an average of approximately 26 words.
%ij: the number of turns seem very low? if its 20 turns per story, shouldn't it be 38*20=760? Or is the number actually the number of interactions, with some participants dropping out?
%yes, it should be the number of interactions. For human-human data, some rows are filtered out for the computational analysis, as they contained nonsense or incomprehensive inputs. Have fixed now.

\subsection{LLM-LLM}
In the LLM-LLM, condition, we replicated the exact constraints and flow from the original platform into a simulation pipeline. For each LLM used in the Human-LLM experiment, the simulation replaced the human with an identical LLM counterpart (i.e. Claude wrote a story with Claude), effectively resulting in a structurally comparable LLM-LLM dataset. 

The condition resulted in 80 stories, each consisting of 10 interactions and an average of approximately 28 words per turn.

\section{Methods}
%Yb: perhaps a table here, for the metrics: metric, unit, tool, and subsection. 




We investigate co-writing dynamics through three dimensions: affective alignment, semantic novelty and surface narrative influence (captured via transience and resonance).

We use a turn-level analytical framework that integrates embedding-based measures of semantic novelty and directional influence with sentiment-based measures of affective alignment.

\subsection{Valence}
\subsection{Directional Alignment}



\subsection{Novelty, transience and resonance}
For narrative influence, we use information-theoretic measures of discourse influence, namely Novelty, Transience and Resonance. Novelty measures the divergence of a turn from prior context, representing the introduction of novel/surprising material to the story. Transience measures the predictive effect of a turn on the immediately following story contribution. Resonance measures the relationship between the two, expressing whether that departure remains predictive of what follows.

\textbf{Novelty}: which actor introduces new material?
\textbf{Transience}: how much does new material fail to persist?
\textbf{Resonance}: whose novelty matters for the future of the story?
{Surprisal}
%Single token surprisal 
We build on the established theory of surprisal-based novelty, transience, and resonance from \cite{fundal2026directionalalignmentnarrativeagency}, and expand the methodological paradigm of transience specifically. The single surprisal score of a token $w_j$ is computed as the negative log-probability of $w_j$ given the preceding tokens in a defined context window $w_{<j}$. The scores were computed using the open-source base model google/gemma-4-31b.
 \begin{equation}
s(w_j) = -\log_2 p(w_j \mid w_{<j})
\end{equation}

%Turn-level surprisal:
The surprisal score of turn-level inputs $\bar{s}(T_t \mid C_t)$ is then computed as the mean surprisal score of the tokens in a turn $T_t$ . Each token is conditioned on the preceding context $C_t$ and the preceding tokens within the turn $w_{t,<j}$.

\begin{equation}
\bar{s}(T_t \mid C_t)
=
\frac{1}{n}
\sum_{j=1}^{n}
-\log_2 p(w_{t,j} \mid C_t, w_{t,<j})
\end{equation}
%Novelty
We then define Novelty and Transience by asking whether the prior context makes the current or future turns less surprising.

Novelty is defined as the difference between contextual and baseline surprisal. This can be interpreted as the negative of a causal model-based PMI estimate between the preceding context and the current turn. The contextual score is computed using all preceding tokens in the story $w_{<t}$. The baseline score is calculated as the surprisal score of the turn from the model's beginning-of-sequence token. Formally: 
\begin{equation}
\mathrm{Novelty}_t = \bar{s}(T_t \mid C_t) - \bar{s}(T_t \mid \mathrm{BOS})
\end{equation}
%Transience
We expand on the definition of transience from \cite{fundal2026directionalalignmentnarrativeagency}. Instead of regarding transience only as the isolated total predictive utility of a turn $s_t$, regardless of the prior context, we now measure the marginal contribution of $s_t$ on predicting the future. This is obtained by conditioning on the past and asking how much the current turn helps predicting the future turn \textbf{given} what has already been established. Formally, we transition from 

$Transience_t=\bar{s}(F_{t+1}\mid T_t)-\bar{s}(F_{t+1}\mid BOS)$

To

$Transience_t=\bar{s}(F_t\mid C_t,T_t)-\bar{s}(F_t\mid C_t)$

Where $F_t$ is the full subsequent counterpart's turn.
This facilitates a measure of what the turn carries into the future narrative, besides what has already been established. 

%Resoncance
Finally, we define resonance as the net gain of novelty on future discourse (rephrase)
\begin{equation}
\mathrm{Resonance}_t = \mathrm{Novelty}_t - \mathrm{Transience}_t
\end{equation}


\subsection{Surface textual measures? Word count, readability, diversity}

%Yb: results table here: alginment, novelty etc. for the three settings

\begin{table*}[t]
\centering
\small
\setlength{\tabcolsep}{4pt}
\renewcommand{\arraystretch}{1.12}
\begin{tabularx}{\textwidth}{>{\RaggedRight\arraybackslash}p{0.18\textwidth}>{\RaggedRight\arraybackslash}p{0.24\textwidth}Y>{\RaggedRight\arraybackslash}p{0.16\textwidth}}
\toprule
Question & Measure & Finding & Evidence \\
\midrule
RQ1: Emotional alignment & Per-story valence alignment asymmetry, $A_{\mathrm{val}}$ & Human--AI showed the largest descriptive asymmetry, but the cross-condition effect was not reliable. Signed Human--AI directionality was suggestive. & $p=.101$; signed HA $p=.023$ \\
\addlinespace
RQ2: Novelty and resonance & Metric-level slot asymmetry, $A_m$ & Human--AI showed substantially larger side-level asymmetry in surprise, transience, and resonance than both baselines. & all planned contrasts $p_{\mathrm{BH}}<.001$ \\
\addlinespace
RQ2: Narrative influence & Novelty-to-resonance slope asymmetry, $A_{\mathrm{infl}}$ & Human--AI showed the largest descriptive slope asymmetry, but the contrast was marginal. & $p=.057$ \\
\addlinespace
RQ3: Interactional regime & Cross-condition comparison across valence and information-theoretic measures & Human--AI appears distinct primarily in metric-level asymmetry and novelty-to-resonance coupling, while valence asymmetry and slope-asymmetry evidence are weaker. & mixed evidence \\
\bottomrule
\end{tabularx}
\caption{Summary of cross-condition findings across the main research questions.}
\label{tab:results}
\end{table*}
\section{Results}
\subsection{Emotional Alignment}
Human-AI collaboration exhibits unique asymmetries: LLMs accommodate human emotional tone. Human-Human and LLM-LLM pairs show symmetric adaptation. 
% Yb: Picture: three panels, each show responder valence predicted by previous/current partner valence, with separate slopes for direction? 
\begin{table}[t]
\centering
\small
\setlength{\tabcolsep}{3pt}
\begin{tabularx}{\columnwidth}{Yrrrr}
\toprule
Condition & $n$ & Mean $A_{\mathrm{val}}$ & 95\% CI & Median \\
\midrule
Human--AI & 96 & 0.606 & [0.517, 0.703] & 0.538 \\
Human--Human & 36 & 0.516 & [0.387, 0.647] & 0.439 \\
AI--AI & 80 & 0.490 & [0.414, 0.570] & 0.425 \\
\bottomrule
\end{tabularx}
\caption{Per-story valence alignment asymmetry across conditions. $A_{\mathrm{val}}$ is the absolute difference between Fisher-transformed directional valence correlations. Human--AI shows the highest descriptive asymmetry, but the omnibus and planned contrast tests are not statistically significant.}
\label{tab:valence-asymmetry}
\end{table}
\subsection{Narrative Agency}
Humans are the main source of narrative novelty and resonance in Human-Ai writing. LLM-LLM collaborations are, in a relative sense, predictable.
% Yb: Image: regression lines? 
The asymmetry in narrative influence is significantly higher for human-LLM stories, than human-human and LLM-LLM stories. 
\begin{table*}[t]
\centering
\small
\setlength{\tabcolsep}{4pt}
\begin{tabular}{llrrcc}
\toprule
Metric & Condition & $n$ & Mean $A_m$ & 95\% CI & Cliff's $\delta$ vs. HA \\
\midrule
Surprise & Human--AI & 97 & 0.828 & [0.732, 0.926] & -- \\
Surprise & Human--Human & 36 & 0.355 & [0.284, 0.428] & 0.591 \\
Surprise & AI--AI & 80 & 0.209 & [0.175, 0.245] & 0.766 \\
\addlinespace
Transience & Human--AI & 97 & 0.484 & [0.432, 0.537] & -- \\
Transience & Human--Human & 36 & 0.300 & [0.233, 0.370] & 0.397 \\
Transience & AI--AI & 80 & 0.197 & [0.165, 0.231] & 0.670 \\
\addlinespace
Resonance & Human--AI & 97 & 1.290 & [1.149, 1.432] & -- \\
Resonance & Human--Human & 36 & 0.605 & [0.482, 0.733] & 0.572 \\
Resonance & AI--AI & 80 & 0.313 & [0.254, 0.382] & 0.806 \\
\bottomrule
\end{tabular}
\caption{Metric-level slot asymmetry across conditions. $A_m$ denotes the per-story absolute difference between slot-level mean metric values. Human--AI shows substantially larger side-level asymmetry in surprise, transience, and resonance than both comparison conditions. Cliff's $\delta$ values compare Human--AI against each baseline condition.}
\label{tab:metric-asymmetry}
\end{table*}
\subsection{Rubber-Band Effect}
Across conditions emotional divergence tends to trigger subsequent convergence, suggesting a general corrective dynamic in collaborative storytelling. 

\subsection{Style and Surface Properties}
Human-Human stories are more readable than stories involving AI. LLM participation leaves a detectable stylistic signature in mixed collaboration, leading to a surface-complex language that is more complex than the human average.

\section{Discussion}
Treating co-writing as an unfolding interaction, our work has provided a finer grained account of creativity and coordination in open ended text generation. We have revealed systematic patterns in how narrative influence and alingment unfold across conditions. 

Together these findings suggest that Human-AI co-writing is a distinct interactional regime characterized by asymmetric agency and adaptive alignment, with properties the differ from both Human-Human and LLM-LLM co-writing. 

Our approach improves over prior evaluations by enabling a fine-grained analysis of co-writing dynamics and capturing both narrative innovation and adaptive behavior. 

Our approach contributes to NLP research on human-AI collaboration by modeling narrative development as an emergence distribution of agency across interactive writers, as well as to creative co-writing research in general ...

\section{Conclusion and Future Works}

\section*{Limitations}

This document does not cover the content requirements for ACL or any
other specific venue.  Check the author instructions for
information on
maximum page lengths, the required ``Limitations'' section,
and so on.

\section*{Acknowledgments}

This document has been adapted
by Steven Bethard, Ryan Cotterell and Rui Yan
from the instructions for earlier ACL and NAACL proceedings, including those for
ACL 2019 by Douwe Kiela and Ivan Vuli\'{c},
NAACL 2019 by Stephanie Lukin and Alla Roskovskaya,
ACL 2018 by Shay Cohen, Kevin Gimpel, and Wei Lu,
NAACL 2018 by Margaret Mitchell and Stephanie Lukin,
Bib\TeX{} suggestions for (NA)ACL 2017/2018 from Jason Eisner,
ACL 2017 by Dan Gildea and Min-Yen Kan,
NAACL 2017 by Margaret Mitchell,
ACL 2012 by Maggie Li and Michael White,
ACL 2010 by Jing-Shin Chang and Philipp Koehn,
ACL 2008 by Johanna D. Moore, Simone Teufel, James Allan, and Sadaoki Furui,
ACL 2005 by Hwee Tou Ng and Kemal Oflazer,
ACL 2002 by Eugene Charniak and Dekang Lin,
and earlier ACL and EACL formats written by several people, including
John Chen, Henry S. Thompson and Donald Walker.
Additional elements were taken from the formatting instructions of the \emph{International Joint Conference on Artificial Intelligence} and the \emph{Conference on Computer Vision and Pattern Recognition}.

% Bibliography entries for the entire Anthology, followed by custom entries
%\bibliography{anthology,custom}
% Custom bibliography entries only
\bibliography{custom}

\appendix

\section{Example Appendix}
\label{sec:appendix}

This is an appendix.

\end{document}
